%! Introduction to Functions and Change — Unit 1, Lesson 1.0 — AP Practice (ANSWER KEY)
% EXTRA CREDIT — optional, exactly TWO pages, placed between the notes and the graded homework.
% Shaped like the AP Precalculus exam: page 1 is Section I, five multiple-choice items with four
% options (A)–(D); page 2 is Section II, one multi-part free-response question on a table with
% interpret-in-context parts. Contexts (cities and ZIP codes, a hiker, a phone battery, a food
% truck) are used nowhere else in the lesson. No function notation — it arrives in 1.1.
% Distractors are the lesson's errors on purpose: MC 2 calls a constant output "not a function",
% MC 3 counts arrivals at an output instead of arrows leaving an input, MC 5 swaps input and
% output. FRQ (d) is the reversal from notes problem 19.
%
% PAGE LOCKSTEP with the other half of the pair: the key marks the correct option in its
% LABEL (no text added to the line) and prose answers are \writespace{H}{…}, fixed height both
% ways. The explicit \newpage fixes the page break in both files.
\documentclass[10pt]{article}
\usepackage{precalculus-article}
\usepackage{precalculus-key}
\newcommand{\TallMath}[1]{$\displaystyle #1\rule[-1.4em]{0pt}{3.2em}$}

\begin{document}

\pageheader{Unit 1, Lesson 1.0}{AP Practice --- Answer Key}

\vspace{0.12in}

\boxguard[8]
\begin{remindbox}
\small
\textbf{Extra credit --- optional.} These two pages are written in the style of the AP
Precalculus exam. Your graded homework is the last section of this packet; do it first. Every
answer choice below is a mistake someone really makes, so read all four before you choose.
\end{remindbox}

\vspace{0.12in}

\begin{headlinebox}{goldbg}
\small
\textbf{Section I --- Multiple Choice.} Circle the letter of the single best answer. No calculator.
\end{headlinebox}

\vspace{0.06in}

\begin{enumerate}[leftmargin=1.9em, itemsep=12pt]

\item Which of the following rules is \textbf{not} a function?
      \begin{enumerate}[leftmargin=2.6em, itemsep=2pt]
        \item[(A)] Each person $\to$ the date they were born
        \item[(B)] Each car $\to$ its vehicle identification number
        \item[\textcolor{keyred}{\textbf{$\checkmark$\,(C)}}] Each city $\to$ a ZIP code in that city
        \item[(D)] Each whole number $\to$ that number plus 10
      \end{enumerate}

\item A parking lot charges the same fee no matter how long a car stays. The table shows the
      fee for several stays.

      \begin{center}
      \renewcommand{\arraystretch}{1.3}
      \begin{tabular}{|c|c|c|c|c|}
      \hline
      Hours parked & 1 & 2 & 3 & 4 \\ \hline
      Fee (\$)     & 6 & 6 & 6 & 6 \\ \hline
      \end{tabular}
      \end{center}

      Which statement is true?
      \begin{enumerate}[leftmargin=2.6em, itemsep=2pt]
        \item[(A)] Fee is not a function of hours parked, because every input has the same output.
        \item[\textcolor{keyred}{\textbf{$\checkmark$\,(B)}}] Fee is a function of hours parked, and the fee is constant.
        \item[(C)] Fee is a function of hours parked, and the fee is increasing.
        \item[(D)] Hours parked is a function of the fee, because the fee is always \$6.
      \end{enumerate}

\item Each list below gives ordered pairs (input, output). Which list does \textbf{not} describe a
      function?

      \smallskip
      \textbf{(A)} $(1, 2),\ (2, 2),\ (3, 2)$ \hfill \textbf{(B)} $(0, 1),\ (1, 0),\ (2, 1)$ \hfill
      \textcolor{keyred}{\textbf{(C)}} $(3, 4),\ (5, 6),\ (3, 8)$ \hfill \textbf{(D)} $(-1, 1),\ (1, 1),\ (0, 0)$

\item \begin{minipage}[t]{0.52\linewidth}
      A hiker's elevation, in hundreds of feet, is graphed against hours on the trail. On which
      interval is the hiker's elevation \textbf{decreasing}?
      \begin{enumerate}[leftmargin=2.6em, itemsep=3pt]
        \item[(A)] From hour 0 to hour 3
        \item[(B)] From hour 3 to hour 5
        \item[\textcolor{keyred}{\textbf{$\checkmark$\,(C)}}] From hour 5 to hour 8
        \item[(D)] From hour 8 to hour 10
      \end{enumerate}
      \end{minipage}\hfill
      \begin{minipage}[t]{0.44\linewidth}
      \vspace{-6pt}
      \centering
      \begin{tikzpicture}[x=0.4cm, y=0.3cm, baseline=(current bounding box.north)]
        \draw[linegray, very thin, step=1] (0,0) grid (10,9);
        \draw[->, thick] (0,0) -- (10.6,0) node[right] {\scriptsize hours};
        \draw[->, thick] (0,0) -- (0,9.6) node[above] {\scriptsize elevation};
        \foreach \x in {2,4,6,8,10} \node[below, font=\tiny] at (\x,0) {\x};
        \foreach \y in {2,4,6,8} \node[left, font=\tiny] at (0,\y) {\y};
        \draw[very thick, plum] (0,2) -- (3,8) -- (5,8) -- (8,4) -- (10,6);
      \end{tikzpicture}
      \end{minipage}

\item A student records how many hours her phone screen has been on and the percent of battery
      remaining. Which statement best describes the relationship?
      \begin{enumerate}[leftmargin=2.6em, itemsep=2pt]
        \item[\textcolor{keyred}{\textbf{$\checkmark$\,(A)}}] Hours on is the input; battery percent is the output, because it responds to the hours.
        \item[(B)] Battery percent is the input; hours on is the output, because the battery runs the screen.
        \item[(C)] Both quantities are inputs, because the student records both of them.
        \item[(D)] Neither is an output, because battery percent goes down instead of up.
      \end{enumerate}

\end{enumerate}

\newpage

\begin{headlinebox}{goldbg}
\small
\textbf{Section II --- Free Response.} Show your reasoning. Answers about the food truck must be
written \emph{in context}, with units --- that is what earns credit on the AP exam.
\end{headlinebox}

\vspace{0.08in}

\begin{enumerate}[leftmargin=1.9em, itemsep=10pt, start=6]

\item A food truck keeps a running count of the total number of tacos it has sold since opening.

      \smallskip
      \begin{center}
      \renewcommand{\arraystretch}{1.35}
      \begin{tabular}{|c|c|c|c|c|c|c|}
      \hline
      Hours since opening & 0 & 1  & 2  & 3  & 4   & 5 \\ \hline
      Total tacos sold    & 0 & 40 & 90 & 90 & 130 & 150 \\ \hline
      \end{tabular}
      \end{center}

      \textbf{(a)} Identify the input and the output, with units. Which is the dependent variable?
      \writespace{2.0cm}{Input: hours since opening (hours). Output: total tacos sold (tacos). The dependent variable is total tacos sold --- it responds to the time.}

      \textbf{(b)} Over which interval shown is the total \textbf{constant}? Describe what could
      have happened at the truck during that interval.
      \writespace{2.0cm}{From hour 2 to hour 3. The total did not change, so no tacos were sold --- for example, the truck closed for a break or ran out of tortillas.}

      \textbf{(c)} Is the total number of tacos sold a function of hours since opening? Explain.
      \writespace{2.0cm}{Yes. At each time there is exactly one running total. Hours 2 and 3 share the total 90, but many-to-one is allowed.}

      \textbf{(d)} The manager says, ``If you tell me the total was 90 tacos, I can tell you exactly
      what hour it was.'' Is hours since opening a function of total tacos sold? Explain.
      \writespace{2.4cm}{No. The input 90 tacos would have two outputs, hour 2 and hour 3. One input with two outputs is not a function, so the manager cannot name one exact hour.}

      \textbf{(e)} Could the total number of tacos sold ever be \textbf{decreasing}? Explain in the
      context of the problem.
      \writespace{2.4cm}{No. It is a running total of tacos already sold; a sale cannot be un-sold. The total can increase or stay constant, but it can never go down.}

\end{enumerate}

\end{document}
